\begin{figure}[ht]
\centering
\begin{tikzpicture}[x=0.6cm, y=0.6cm, font=\footnotesize]

% Two side-by-side panels: well-conditioned vs ill-conditioned 2x2 system
% Each panel shows two lines (rows of A x = b) and their intersection;
% small perturbation of b becomes small/large displacement of solution.

% --- Left panel: well-conditioned
\begin{scope}[shift={(0,0)}]
  \node[anchor=south] at (3,4.4) {\textbf{Well-conditioned} $\kappa \approx 1$};
  \draw[->, thick] (-0.2,0) -- (6,0) node[right] {$x_{1}$};
  \draw[->, thick] (0,-0.2) -- (0,4) node[above] {$x_{2}$};

  % Line 1: x1 + x2 = 4
  \draw[thick, engblue] (0,4) -- (4,0);
  \node[engblue, anchor=south west, font=\scriptsize] at (4.0,0.0) {$L_{1}$};
  % Line 2: x1 - x2 = 0, slope 1
  \draw[thick, engred] (0,0) -- (3.5,3.5);
  \node[engred, anchor=south, font=\scriptsize] at (3.6,3.5) {$L_{2}$};
  % Perturbed line 2: shifted by 0.2 down
  \draw[thick, engred, dashed] (0.2,0) -- (3.7,3.5);

  % Solutions: intersection at (2,2) and ~ (2.1, 1.9)
  \fill[black] (2,2) circle (2pt);
  \fill[black] (2.1,1.9) circle (2pt);
  \draw[thick, ->, gray] (2,2) -- (2.1,1.9);
  \node[anchor=west, font=\scriptsize] at (2.2,1.85)
    {small $\Delta x$};
\end{scope}

% --- Right panel: ill-conditioned (near-parallel lines)
\begin{scope}[shift={(8,0)}]
  \node[anchor=south] at (3,4.4) {\textbf{Ill-conditioned} $\kappa \gg 1$};
  \draw[->, thick] (-0.2,0) -- (6,0) node[right] {$x_{1}$};
  \draw[->, thick] (0,-0.2) -- (0,4) node[above] {$x_{2}$};

  % Line 1: gentle slope
  \draw[thick, engblue] (0,0.5) -- (5.5,3.25);
  \node[engblue, anchor=south, font=\scriptsize] at (5.5,3.25) {$L_{1}$};
  % Line 2: slightly different slope (near-parallel)
  \draw[thick, engred] (0,0.2) -- (5.5,3.5);
  \node[engred, anchor=west, font=\scriptsize] at (5.5,3.5) {$L_{2}$};
  % Perturbed line 2: shift by 0.1
  \draw[thick, engred, dashed] (0,0.4) -- (5.5,3.7);

  % Intersections: very far apart
  \fill[black] (2.4,1.7) circle (2pt);
  \fill[black] (4.6,2.8) circle (2pt);
  \draw[thick, ->, gray] (2.4,1.7) -- (4.55,2.78);
  \node[anchor=south west, font=\scriptsize] at (3.2,2.2)
    {large $\Delta x$};
\end{scope}

\end{tikzpicture}
\caption[Geometric reading of the condition number]{Geometric reading of
the condition number for a $2 \times 2$ linear system. The two rows of
$\mat{A}$ are two lines in the $(x_{1}, x_{2})$ plane; the solution
$\vect{x}$ is their intersection. When the lines are nearly orthogonal
the system is well-conditioned: a small perturbation of $\vect{b}$
(shifting one of the lines by a small amount) moves the intersection by
a comparable small amount. When the lines are nearly parallel the
system is ill-conditioned: the same shift moves the intersection by a
large amount, amplified by $\kappa(\mat{A})$. The Hilbert matrix
exercise of section 16.4 makes this concrete in $n$ dimensions.
Source: authors' diagram; geometric framing follows
\cite[ch.~12]{text:trefethen-bau1997}.}
\label{fig:vol02:ch16:condition-number}
\end{figure}
